\documentclass{report}
\usepackage{amsmath,amssymb}
\setlength{\parindent}{0mm}
\setlength{\parskip}{1em}
\begin{document}
\begin{center}
\rule{6in}{1pt} \
{\large Michael J. Brazell \\
{\bf Scalable Non-linear Solution Techniques for Higher-Order Accurate Finite-Element Reynolds-Averaged Navier-Stokes Problems}}

Department of Mechanical Engineering \\ University of Wyoming \\ Laramie \\ WY 82071 \\ USA
\\
{\tt mbrazell@uwyo.edu}\\
Dimitri J. Mavriplis\\
Behzad R. Ahrabi\end{center}

Although second-order accurate finite-volume methods for structured and
unstructured grids are well established as the dominant approach for
solving the Reynolds-averaged Navier-Stokes (RANS) equations for the
simulation of turbulent flow problems, interest in higher-order
finite-element discretizations has been growing over the last decade.
Considered methods include both continuous and discontinuous Galerkin
discretizations, with the continuous methods being based on the
streamwise-upwind Petrov Galerkin approach (SUPG). The advantages of
these discretizations include the use of a nearest neighbor stencil,
straight forward extension to high-order accuracy, and the generation of
dense computations kernels that are more suitable for emerging exascale
architectures. However, the resulting non-linear equations that must be
solved for steady-state and time-implicit RANS problems have proven to be
much stiffer and more difficult to converge efficiently and robustly.

In this work we examine several non-linear solution techniques for SUPG
and DG discretizations of the RANS equations with a single equation
turbulence model. The approach relies on a preconditioned Newton-Krylov
method, where the exact Jacobian of the discretization is assembled and
used in the GMRES phase of the Newton Krylov solver. Incomplete
factorization (ILU) with various levels of fill is used as a
preconditioner, as well as simpler and more scalable approaches such as
line preconditioning where lines are extracted from the Jacobian matrix
using a graph algorithm. Non-linear continuation is incorporated into the
solver by controlling the pseudo time step applied to the diagonal of the
Jacobian matrix, while a line search is used to ensure suitable
non-linear updates for the Newton solver. We compare the performance and
scalability of the current approach with the performance of similar and
other multigrid solvers applied to traditional second-order accurate
finite-volume unstructured mesh RANS discretizations, and discuss
possible future avenues for improving the efficiency and scalability of
higher-order RANS solvers.


\end{document}
